% fixture: setspace/mixed
% tier: T3
% engine: document
% covers: setspace's \setstretch{1.3} at 12pt with hyperref (anchors at the stretched \baselineskip): a tabular in text (\@arstrutbox from the stretched \strutbox), a table float single-spaced, a minipage, an enumerate, a short display (\everydisplay stretches the short skips too), a footnote
% issue: #253
\documentclass[12pt]{article}
\usepackage{setspace}
\usepackage{hyperref}
\setstretch{1.3}
\begin{document}
\section{Tables}\label{s}
Line spacing changes how many lines fit on a page, where the page breaks fall, and how far apart the baselines of a paragraph sit. A tabular \begin{tabular}{|l|c|}\hline a & b \\ c & d \\\hline\end{tabular} in the text.\footnote{A footnote with a reference to Section~\ref{s}, long enough to run over two lines.}

\begin{table}[h]
\centering
\begin{tabular}{lr}
one & 1 \\
two & 2 \\
\end{tabular}
\caption{A table whose caption is long enough to need two lines in the float, single-spaced there.}\label{t}
\end{table}
Line spacing changes how many lines fit on a page, where the page breaks fall, and how far apart the baselines of a paragraph sit. See Table~\ref{t}.
\begin{enumerate}
\item First item, which runs on for long enough to need a second line in the list.
\item Second.
\end{enumerate}
\begin{minipage}{0.6\linewidth}
Line spacing changes how many lines fit on a page, where the page breaks fall, and how far apart the baselines of a paragraph sit.
\end{minipage}

Short
\[ a = b \]
Line spacing changes how many lines fit on a page, where the page breaks fall, and how far apart the baselines of a paragraph sit.
\end{document}
